\section{Method}

The client selects literal project dependencies and sends them to a private
TeX server. \Cref{fig:pipeline} is an external PNG file, so its presence checks
both dependency discovery and binary artifact transport.

\begin{figure}[H]
  \centering
  \includegraphics[width=\columnwidth]{remote-compilation}
  \caption{The experimental apparatus. Sources go right; the PDF comes back.
  This directional insight required a diagram.}
  \label{fig:pipeline}
\end{figure}

Let each indicator below be one when the corresponding dependency is present
and zero otherwise:

\begin{equation}
  S = I_{\mathrm{class}} I_{\mathrm{input}} I_{\mathrm{figure}}
      I_{\mathrm{bib}} I_{\mathrm{engine}}.
  \label{eq:success}
\end{equation}

\Cref{eq:success} predicts success only when every required part is
available. Multiplying several ones still gives one, which is reassuring and
also the main mathematical contribution of this paper.

\begin{table}[H]
  \caption{Dependencies and the evidence that they were not forgotten}
  \label{tab:dependencies}
  \centering
  \begin{tabular}{@{}ll@{}}
    \toprule
    Dependency & Observable evidence \\
    \midrule
    Included sections & This sentence exists \\
    External figure & \Cref{fig:pipeline} is visible \\
    Equation labels & \Cref{eq:success} has a number \\
    BibTeX database & Citations appear in the references \\
    \bottomrule
  \end{tabular}
\end{table}

Listing~\ref{lst:commands} is loaded from another file. It records the procedure
without asking TeX to execute any command.

\lstinputlisting[
  caption={Printed compile procedure.},
  captionpos=b,
  label={lst:commands},
  linewidth=0.96\columnwidth,
  xleftmargin=0.02\columnwidth
]{code/procedure.txt}
